\subsection{A bounded refinement for \texorpdfstring{$R_{1154}$}{R1154}}
\label{subsec:followup-r1154}

We apply Lemma~\ref{lem:followup-pair-gram} to the complete domain of
$1,238$ cubic-surviving integral Gram patterns at low index $31$.
These patterns are indexed in the frozen $14,211$-pattern integral Gram file;
there are $16$ ordered high restriction types. The exact verifier rebuilds the
cubic coefficient rows and the complete survivor list, checks all $35$ scalar
fourth moments, and verifies the ordered high restriction domain.
All patterns have the same matrix $Q=MM^{\mathsf T}$ of rank $3$.
The simultaneous pair-product Gram matrix $G$ also has rank $3$.

Let $X_0$ be the minimum-norm solution of $MX=T$. In the pair order
$(11,12,13,14,22,23,24,33,34,44)$, put
\[
 h=(1,0,1,-1,0,0,0,1,-1,1).
\]
For every pattern the exact residual is
\[
 G-X_0^{\mathsf T}X_0=\alpha h^{\mathsf T}h,
 \qquad
 \alpha\in\{92/205,262/205,672/205,1082/205\}.
\]
In particular, all $1,238$ patterns pass positivity and the residual rank
bound. Every required nonnegative integral factorization must therefore have
\[
 X=X_0+t h,\qquad Mt=0,\qquad \sum_z t_z^2=\alpha.
\]
Indeed, the minimum-norm solution is orthogonal to $\ker M$, so subtracting
its Gram leaves the Gram of $X-X_0$; the displayed rank-one form forces each
row of this difference to be a scalar multiple of $h$.

For an individual high row, set $q=(X_0)_{z,11}$. Since $h_{11}=1$,
integrality forces $t_z=n-q$ with $n\in\mathbb Z$. The other entries give
integer congruences and nonnegativity bounds $L\leq n\leq U$, all checked
exactly. The minimum of $(n-q)^2$ on this integer interval occurs at the
clipped floor or ceiling of $q$. Summing these minima with the exact high
multiplicities gives a necessary lower bound for $\alpha$.
This excludes $892$ patterns.

A second screen retains the finite individual domains satisfying
$(n-q)^2\leq\alpha$. It uses lower and upper bounds for the sum of squares
and coordinatewise interval bounds for $Mt=0$. A candidate value is removed
when fixing that one row makes a lower norm bound or a linear interval
inconsistent; the same removal is valid for all interchangeable rows of that
restriction type. Iterating these necessary tests requires at most three
passes, with at most four candidate values per coordinate. No coupled
integer assignment is enumerated. A further $54$ patterns are excluded:
$34$ by an upper norm bound, $16$ by a linear interval bound, and $4$ by a
lower norm bound.
\[
\begin{array}{c|r|r|r|r}
\alpha&\text{patterns}&\text{minimum-square exclusions}&
\text{further exclusions}&\text{remaining}\\\hline
92/205&8&8&0&0\\
262/205&864&852&6&6\\
672/205&300&32&48&220\\
1082/205&66&0&0&66\\\hline
\text{total}&1238&892&54&292
\end{array}
\]

The eight exclusions in the frozen height-six campaign have indices
$207,208,376,377,898,899,1407,1408$ and are exactly the first row of the table.
Thus the combined exclusion set has size $946$, including $938$ additional
patterns beyond that campaign; all $292$ surviving domains also survived
height six. This is a refinement of necessary central data, not a fusion-ring
exclusion or a construction of a center. Joint integral solvability of
$Mt=0$ and the exact norm equation remains untested on those $292$ domains.
The general-complex scope uses only the previously accepted
Galois-to-pseudo-unitary reduction.

The portable replay scripts \texttt{r1154\_quartic\_necessary\_filters.py},
\texttt{r1154\_quartic\_coordinate\_propagation.py}, and
\texttt{r1154\_quartic\_reconcile.py} bind their inputs by SHA-256 and give
byte-identical normal and optimized-Python receipts. The last script only
reconciles the already verified height-six certificate indices; it does not
rerun that campaign.
